%%%%%%%%%%%%%%%%%%%%%%%%%%%%%%%%%%%%%%%%%
% Lachaise Assignment
% Structure Specification File
% Version 1.0 (26/6/2018)
%
% This template originates from:
% http://www.LaTeXTemplates.com
%
% Authors:
% Marion Lachaise & François Févotte
% Vel (vel@LaTeXTemplates.com)
%
% License:
% CC BY-NC-SA 3.0 (http://creativecommons.org/licenses/by-nc-sa/3.0/)
% 
%%%%%%%%%%%%%%%%%%%%%%%%%%%%%%%%%%%%%%%%%

%----------------------------------------------------------------------------------------
%	PACKAGES AND OTHER DOCUMENT CONFIGURATIONS
%----------------------------------------------------------------------------------------
\usepackage{amsmath,amsthm,amsfonts,amssymb,graphicx,mathrsfs,mathtools,bm,stmaryrd,xcolor,enumerate,physics,tikz}
\usepackage[margin=1in]{geometry}


%----------------------------------------------------------------------------------------
%	FONTS
%----------------------------------------------------------------------------------------

\usepackage[utf8]{inputenc} % Required for inputting international characters
\usepackage[T1]{fontenc} % Output font encoding for international characters

\newcommand{\NN}{\mathbb{N}}
\newcommand{\ZZ}{\mathbb{Z}}
\newcommand{\QQ}{\mathbb{Q}}
\newcommand{\RR}{\mathbb{R}}
\newcommand{\CC}{\mathbb{C}}
\newcommand{\HH}{\mathbb{H}}
\newcommand{\OO}{\mathbb{O}}
\newcommand{\FF}{\mathbb{F}}

% Script Characters
\newcommand{\T}{\mathcal{T}} %mathscr also exists, but it tends to be too fancy for anything other than algebraic geometry
\newcommand{\B}{\mathcal{B}}
\newcommand{\D}{\mathcal{D}}
\newcommand{\R}{\mathcal{R}}
\newcommand{\K}{\mathcal{K}}
\newcommand{\N}{\mathcal{N}}
\newcommand{\Z}{\mathcal{Z}}
\newcommand{\Q}{\mathcal{Q}}

% Boldscript Numbers
\newcommand{\0}{\mathbb{0}}
\newcommand{\1}{\mathbb{1}}

% Functions
\newcommand{\PP}{\mathbb{P}}
\newcommand{\EE}{\mathbb{E}}

% New commands
\newcommand{\ra}{\rightarrow}
% \newcommand{\la}{\leftarrow}
\newcommand{\nor}{\vartriangleleft}
\newcommand{\bracket}[1]{[#1]}
\newcommand{\aut}{\text{Aut }}
\newcommand{\spec}{\text{Spec }}
\newcommand{\inn}{\text{Inn}}
\newcommand{\solidline}{\\\noindent\makebox[\linewidth]{\rule{\paperwidth}{0.4pt}}\\}


% Renewing commands
\renewcommand{\Im}{\text{Im }}
\renewcommand{\Re}{\text{Re }}
\renewcommand{\hom}{\text{Hom}}
\renewcommand{\bar}{\overline}
\renewcommand{\l}{\langle}
\renewcommand{\r}{\rangle}
\newcommand{\vbrack}[1]{\left\langle #1\right\rangle}
\renewcommand{\theenumi}{\alph{enumi}}
\renewcommand\qedsymbol{$\blacksquare$}

% Griffith's commands
\usetikzlibrary{arrows,scopes} %tikz libraries used to draw the glyphs
%defining script r (scalar)
\newcommand{\rc}{%
\resizebox{!}{1.25ex}{%
    \begin{tikzpicture}[>=round cap]
        \clip (0.09em,-0.05ex) rectangle (0.61em,0.81ex);
        \draw [line width=.11ex, <->, rounded corners=0.13ex] (0.1em,0.1ex) .. controls (0.24em,0.4ex) .. (0.35em,0.8ex) .. controls (0.29em,0.725ex) .. (0.25em,0.6ex) .. controls (0.7em,0.8ex) and (0.08em,-0.4ex) .. (0.55em,0.25ex);
    \end{tikzpicture}%
}%
}

%defining bold script r (vector)
\newcommand{\brc}{%
\resizebox{!}{1.3ex}{%
    \begin{tikzpicture}[>=round cap]
        \clip (0.085em,-0.1ex) rectangle (0.61em,0.875ex);
        \draw [line width=.2ex, <->, rounded corners=0.13ex] (0.1em,0.1ex) .. controls (0.24em,0.4ex) .. (0.35em,0.8ex) .. controls (0.29em,0.725ex) .. (0.25em,0.6ex) .. controls (0.7em,0.8ex) and (0.08em,-0.4ex) .. (0.55em,0.25ex);
    \end{tikzpicture}%
}%
}

%defining bold script r with a hat (unit vector)
\newcommand{\hrc}{\hat{\brc}}

\newcommand{\perm}{4 \pi \epsilon_0}
\newcommand{\bvec}[1]{\hat{\textbf{#1}}}